\documentclass[mlmain]{jmlr}
\jmlrpmlr

% Link to code: https://github.com/lanl/DeepLean

% Preserve ligatures and accented glyphs under PDF text extraction.
\usepackage[T1]{fontenc}
\usepackage[protrusion=false,expansion=false]{microtype}
\DisableLigatures{encoding=*,family=*}
\input{glyphtounicode}
\pdfgentounicode=1
\usepackage{listings}

% Reviewer-facing Lean declarations are read directly from the checked source.
\lstset{
  basicstyle=\ttfamily\scriptsize,
  breaklines=true,
  columns=fullflexible,
  keepspaces=true,
  showstringspaces=false,
  frame=single,
  literate=
    {ℝ}{{$\mathbb{R}$}}1
    {ℕ}{{$\mathbb{N}$}}1
    {τ}{{$\tau$}}1
    {∀}{{$\forall$}}1
    {∃}{{$\exists$}}1
    {→}{{$\to$}}1
    {↦}{{$\mapsto$}}1
    {∧}{{$\wedge$}}1
    {¬}{{$\neg$}}1
    {↔}{{$\leftrightarrow$}}1
    {∈}{{$\in$}}1
    {≤}{{$\leq$}}1
    {≠}{{$\neq$}}1
    {⟨}{{$\langle$}}1
    {⟩}{{$\rangle$}}1
    {₁}{{$_1$}}1
    {₂}{{$_2$}}1
}

% Review watermark disabled for the camera-ready version.
%\usepackage[hpos=300px,vpos=70px]{draftwatermark}
%\SetWatermarkText{\test}
%\SetWatermarkScale{1}
%\SetWatermarkAngle{0}

\setlength{\emergencystretch}{2em}

\newcommand{\R}{\mathbb{R}}
\newcommand{\cF}{\mathcal{F}}
\newcommand{\cP}{\mathcal{P}}
\newcommand{\cU}{\mathcal{U}}
\newcommand{\T}{\mathsf{T}}
\newcommand{\eps}{\varepsilon}

% \newtheorem{myremark}{Remark}
\newtheorem{mylemma}{Lemma}

\newcounter{myremark}

\newenvironment{myremark}
  {\refstepcounter{myremark}%
   \par\medskip
   \noindent\textbf{Remark~\themyremark.}\ }
  {\par\medskip}

\jmlrvolume{}
\firstpageno{1}
\editors{}
\jmlryear{2026}
\jmlrworkshop{Symmetry and Geometry in Neural Representations}

\title[Identifiability of Deep Normalized Attention]
{Identifiability of Deep Normalized Attention: Or,\\ Machine Verified Proofs for Neuroalgebraic Geometry
}
% \author{Anonymous Authors}
\author{%
  \Name{James Amarel} \Email{jlamarel@lanl.gov} \\
  \Name{Benjamin Migliori} \Email{ben.migliori@lanl.gov}\\
  \Name{Nathan Debardeleben}  \Email{ndebard@lanl.gov}\\
  \Name{Emily Taylor} \Email{ecasleton@lanl.gov}\\
  \Name{Earl Lawrence} \Email{earl@lanl.gov}\\
  \Name{Gerd J. Kunde} \Email{g.j.kunde@lanl.gov}\\
  \addr Los Alamos National Laboratory, Los Alamos, NM 87545
}

\begin{document}

\maketitle

\begin{abstract}
    Identifiability of deep normalized attention is not ensured by
positivity and injectivity of the scalar normalizer: we construct a positive injective Borel
normalizer for which every point of an explicit open dense set of scalar depth-two parameters has a distinct partner
with exactly the same function realization.  For exponential softmax, however, we prove that the
virtual weights are uniquely determined at every finite depth and every token
count $t\ge2$, provided that the final linear map is nonzero and the score
matrices before the final attention layer are not skew-symmetric.  The
arbitrary-depth proof, developed with Codex and formally verified in Lean~4,
uses inputs with one distinguished token and repeated copies of another.
At the first layer where two networks have different score matrices, suitable
probes produce a nonzero leading term in the displacement between
corresponding output tokens as the distance between the incoming token
classes tends to zero; the outgoing distance within each network is subdominant,
so later normalized layers cannot cancel the leading displacement.

\end{abstract}

\begin{keywords}
parameter identifiability, functional fibers, normalized self-attention
\end{keywords}

\section{Introduction}
\label{sec:introduction}

Henry, Marchetti, and Kohn conjectured that a deep normalized self-attention
network is identifiable from its realized function after passing to the
virtual parametrization they introduced \citep{henry2025}.  Their lightning,
or unnormalized, model retains layerwise rescaling symmetry even after primitive
query-key factorizations and changes of hidden-state coordinates have been
absorbed into the virtual weights; restoring row normalization suggested that
the remaining rescalings should disappear.  The resulting conjecture, and the
numerical evidence that softmax breaks the lightning rescaling symmetry, were later
highlighted in the broader review of neural-network parameter symmetries by
\citet{zhao2025}.  
In the same single-head attention setting, write
$W$ for the virtual weights and let the realization map send $W$ to the function
$\phi_W^S$ that it defines on token matrices, where
$S:\R\to\R_{>0}$ is injective and $S(0)=1$ \citep{henry2025}.  The complete functional fiber of
$W$ is the set of all virtual parameters that identically realize $\phi_W^S$ on the
whole input space.  \citet{henry2025} proved that this
fiber is a singleton at one layer for tokens $t\ge2$ when the value map is nonzero,
conjectured generic singleton fibers at arbitrary depth, and supported the
conjecture by a sampled-Jacobian calculation at depth two.  We
resolve that conjecture in the virtual parametrization in which it was posed:
compositions of single-head attention layers without residual connections,
layer normalization, feed-forward modules, biases, or positional encodings.

Composing normalized attention layers makes the deep problem different from
the one-layer theorem: each row retains only ratios of scalar weights, and a
later score matrix is evaluated only on token configurations produced by the
earlier normalized layers.  We construct a positive injective Borel normalizer
with $S(0)=1$ and $S(4z)=S(z)/4$ for $z\ne0$.  

Consider a scalar two-layer,
two-token network with score parameters $m_1,m_2$ and final linear map $l$,
where $m_1m_2l\ne0$.  The construction ensures that the first layer sends
every nonzero input to two nonzero tokens, so $m_2\ne0$ makes every
second-layer score nonzero.  Replacing $m_2$ by $4m_2$ therefore multiplies
every unnormalized weight in each second-layer row by $1/4$, which cancels
under normalization.  The two networks also agree at zero, so the distinct
tuples $(m_1,m_2,l)$ and $(m_1,4m_2,l)$ realize the same function.  Every
point with $m_1m_2l\ne0$ thus lies in a non-singleton fiber, and these points
form the open dense complement of the three coordinate hyperplanes.

For standard fixed-temperature softmax, the preceding nonidentifiability
cannot occur under the following explicit conditions: the final linear map
is nonzero and each score matrix before the final attention layer is not
skew-symmetric.  At every finite depth and every token count $t\ge2$, the
realized function then uniquely determines all virtual weights.  This
uniqueness excludes every other parameter point of the same dimensions and
depth, without assuming that the other point satisfies the same
nondegeneracy conditions, and the hypotheses hold on an explicit open dense
set.  The proof first determines the final linear map from constant inputs,
then compares two networks on inputs formed from one distinguished token
and repeated copies of another.  At the first differing score matrix, the
two networks receive the same two token classes.  As the distance $s$
between these classes tends to zero, suitable probes produce a nonzero
leading $s^2$ term in the displacement between corresponding output tokens,
while the outgoing distance within each network is $O(s^3)$.  Subsequent
normalized layers average tokens within these smaller gaps and therefore
preserve the leading displacement.  The probes are chosen so that the final
linear map does not annihilate the leading displacement, allowing the score
matrices to be identified successively through the depth.

\section{Related work}
\label{sec:related}
Exact recovery from a realized function has been established for specific
classes of feedforward networks.  Minimal scalar-output networks with one
hidden tanh layer are identifiable up to their finite structural symmetries
\citep{sussmann1992}, while analytic continuation recovers generic fully
connected deep tanh networks \citep{fefferman1994}.
\citet[Theorem~1]{vlacic2022} construct a family of activations for which
networks of arbitrary depth and connectivity are identifiable up to
hidden-neuron permutations under their nondegeneracy and no-clones
hypotheses.  For ReLU networks, \citet{phuong2020} study functional
equivalence modulo positive rescalings and neuron permutations under
architectural and geometric restrictions, while \citet{stock2023} construct
an embedding invariant under positive neuron rescaling.

Neuron permutations and positive rescalings need not exhaust the functional
fibers of ReLU networks \citep{grigsby2023}.  At the infinitesimal level,
functional dimension measures variation of the realized function through
ranks of finite-batch evaluation Jacobians
\citep{grigsby2026functionaldimension,grigsby2024persistent}.
For depth-two softmax networks, the sampled Jacobian ranks reported by
\citet{henry2025} provide local evidence for the normalized attention
conjecture, but even injectivity of the realization map in a neighborhood of
a virtual parameter point does not exclude another point outside that
neighborhood realizing the same function on every input.

\citet{brunner2020} show that distinct attention weights can produce the same
head output for a fixed input and fixed value vectors.  For scalar-output
attention regressors, \citet{bhattamishra2026} recover a single head's score
matrix and nonzero value vector from exact input-output queries
and exhibit multihead nonuniqueness when heads share a score matrix;
\citet{kim2026queries} recover canonical multihead score-matrix/value-vector
pairs up to head permutation using exact variable-length queries.  For
primitive multihead parameters,
\citet{tran2025nfn} characterize functional equivalence under nondegeneracy
conditions over inputs of every sequence length, while
\citet{henry2026multihead} classifies generic fibers at a fixed token count
$t\ge3$.  \citet{tran2026attention} show that additive sinusoidal encodings
preserve functional equivalence, whereas rotary encodings restrict the
query-key symmetry group.\looseness=-1

For a single lightning-attention layer, \citet{alexandr2026} study polynomial
relations among the coefficients of the realized function.  At arbitrary
depth, \citet*{henry2025} pass to virtual weights, removing primitive
factorization and hidden-coordinate ambiguities.  Their lightning fibers
retain layerwise rescalings, and their normalized-attention conjecture asks
whether the entire input-output function generically determines the virtual
weights uniquely.  We resolve this conjecture, constructing non-singleton
fibers for an admissible normalizer and proving generic singleton fibers for
fixed-temperature softmax.

\section{Virtual normalized attention and functional fibers}
\label{sec:model}

Let $d,t\ge1$ be integers denoting the embedding dimension and token count, respectively.
Write $X=(x_1,\ldots,x_t)\in\R^{d\times t}$ for an indexed sequence of
token columns $x_i\in\R^d$.  For a score matrix $M\in\R^{d\times d}$ and
a positive normalizer $S:\R\to\R_{>0}$, define the layer action through
its $i$th output token column,
\begin{equation}
 [\T_M^S(X)]_i
 =\frac{\sum_{j=1}^t S(x_i^\top Mx_j)x_j}
        {\sum_{j=1}^t S(x_i^\top Mx_j)},
 \label{eq:layer}
\end{equation}
where each position $i=1,\ldots,t$ attends to every input token, including
$j>i$, so the model has no temporal or causal masking.
For positive integers $\ell$ and $d_{\mathrm{out}}$, we consider 
depth
$\ell$ and output dimension $d_{\mathrm{out}}$ normalized attention networks realized by the composition 
\begin{equation}
 \phi_W^S(X)=\bigl(L\circ\T_{M_\ell}^S\circ\cdots\circ\T_{M_1}^S\bigr)(X),
 \label{eq:realization}
\end{equation}
where
\begin{equation}
 W=(M_1,\ldots,M_\ell,L)\in
 (\R^{d\times d})^\ell\times\R^{d_{\mathrm{out}}\times d} =
 \cP_{d,d_{\mathrm{out}},\ell}
\end{equation}
is the tuple of virtual weights.  The final map
$L:\R^d\to\R^{d_{\mathrm{out}}}$ acts on each token. For a fixed normalizer
$S$ and token count $t$, we study the complete virtual fiber of $W$,
\begin{equation}
 \cF_S(W)=\{W'\in\cP_{d,d_{\mathrm{out}},\ell}:
                    \phi_{W'}^S=\phi_W^S\},
 \label{eq:fiber}
\end{equation}
consisting of all virtual parameter tuples of the same dimensions and depth
that realize the same function on every input token matrix; identifiability
is precisely the condition $\cF_S(W)=\{W\}$, under which equality of the
realized functions forces equality of every score matrix and the final
linear map.
In the notation of \citet[Section~3.3, Eq.~(8)]{henry2025}, let $V_r$ be the
value map and $A_r$ the primitive score matrix at layer $r$, and set
$L_r=V_r\cdots V_1$.  Then $M_1=A_1$, $M_r=L_{r-1}^\top A_rL_{r-1}$ for
$r>1$, and $L=L_\ell$.  Moving these cumulative value maps through the
normalized layers gives Equation~\eqref{eq:realization}, so Conjecture~3.10
concerns the virtual fiber after the primitive factorization data have been
absorbed.

For a fixed query token $x_i\in\R^d$, let
$s_{ij}=x_i^\top Mx_j$ denote the score assigned to key $x_j$, and let
$p_{ij}$ denote the corresponding normalized attention weight.  Then
\begin{equation}
 \frac{p_{ij}}{p_{ik}}
 =
 \frac{S(s_{ij})}{S(s_{ik})},
 \label{eq:row-odds}
\end{equation}
so normalization retains the ratios among $S(s_{i1}),\ldots,S(s_{it})$, but
discards their common scale: multiplying all these values by the same
positive constant leaves the attention row unchanged.  Injectivity of $S$
permits recovery of $s_{ij}$ from $S(s_{ij})$, but does not by itself permit
recovery of the scores from a normalized attention row.  In contrast, exponential
softmax $S_\tau(z)=e^{z/\tau}$, with fixed $\tau>0$, implies
$\tau\log(p_{ij}/p_{ik})=s_{ij}-s_{ik}$, so the normalized weights determine
the score differences within each row.  For one layer, $t\ge2$, and a
nonzero value map, the realized function determines the score matrix
\citep{henry2025}.  At greater depth, however, each subsequent score matrix
is evaluated on token configurations produced by the preceding normalized
layers; identifiability thus depends on whether, under the successive action
of the normalized layers, distinct sequences of score matrices produce
distinct input-output functions.

\begin{theorem}[Counterexample to Conjecture~3.10 of \citet{henry2025}]
\label{thm:negative}
There exists a positive injective Borel map $S:\R\to\R_{>0}$ with $S(0)=1$ such
that, for $d=d_{\mathrm{out}}=1$, $t=2$, and $\ell=2$, every
$W=(m_1,m_2,l)$ with $m_1m_2l\ne0$ has the distinct partner
$W^\sharp=(m_1,4m_2,l)$ obeying $\phi_{W^\sharp}^S=\phi_W^S$.
Moreover, the parameter set $\{m_1m_2l\ne0\}$ is nonempty, Euclidean open, and dense.
\end{theorem}

\begin{proof}
Row normalization is unchanged when every unnormalized weight in a row is
multiplied by the same positive constant.  For $x\ne0$, let
$q(x)=\log_4|x|$, $k(x)=\lfloor q(x)\rfloor$, and
$r(x)=q(x)-k(x)\in[0,1)$.  Define
\begin{equation}
S(x) = 4^{g(x)}, \text{ where}\quad
 g(x)=
 \begin{cases}
  -k(x)+\frac34+\frac18r(x),&x>0,\\
  -k(x)-\frac34+\frac18r(x),&x<0,\\
  0,&x=0,
 \end{cases}
 \label{eq:exotic-normalizer}
\end{equation}
with $k(x)$ specifying the magnitude interval
$4^{k(x)}\le|x|<4^{k(x)+1}$ and $r(x)$ determining the relative magnitude
through $|x|/4^{k(x)}=4^{r(x)}$.  The fractional part of $g(x)$ lies in
$[3/4,7/8)$ for $x>0$ and $[1/4,3/8)$ for $x<0$, thereby determining the
sign of $x$.  Appendix~\ref{app:counterexample} shows that $x$ can be
recovered from $g(x)$, proving injectivity of $g$ and hence of $S$.  For
$x\ne0$, multiplication by four increases $k(x)$ by one while preserving
$r(x)$ and the sign, so $g(4x)=g(x)-1$ and $S(4x)=S(x)/4$.  Positivity and
$S(0)=1$ follow from the definition; Appendix~\ref{app:counterexample}
verifies Borel measurability.

Let $(a,b)\ne(0,0)$ be an input and $(y_1,y_2)=\T_{m_1}^S(a,b)$ its
first-layer output.  The second-layer scores are $m_2y_i y_j$,
$i,j\in\{1,2\}$, so $m_2\ne0$ reduces their nonvanishing to $y_1y_2\ne0$.
Appendix~\ref{app:counterexample} proves $zS(z)>wS(-w)$ for all $z,w>0$,
which excludes cancellation in either first-layer row when $ab<0$ and
$m_1\ne0$.  If $ab\ge0$, the tokens have the same weak sign and at least
one is nonzero, so every positive weighted average of them is nonzero.
Thus $y_1y_2\ne0$.

Since $W$ and $W^\sharp$ differ only in their second score matrix, they have
the same first-layer output $(y_1,y_2)$.  For every nonzero input,
$m_2y_i y_j\ne0$ for all $i,j$, and hence
$S(4m_2y_i y_j)=S(m_2y_i y_j)/4$.  Replacing $m_2$ by $4m_2$ therefore
multiplies every unnormalized weight in each second-layer row by the same
factor $1/4$, which cancels under row normalization.  The zero input agrees
separately, and the common final map $l$ preserves equality of the two
realizations.  Because $m_2\ne0$, the parameter points are distinct.  Finally,
$\{m_1m_2l\ne0\}$ is the complement of the coordinate hyperplanes $m_1=0$,
$m_2=0$, and $l=0$, and is therefore nonempty, open, and dense.
\end{proof}


\begin{myremark}
For a one-layer scalar network with the normalizer of
Theorem~\ref{thm:negative}, score parameter $m\ne0$, and final map $l\ne0$,
an input $(a,0)$ with $a\ne0$ gives the query $a$ one nonzero score $ma^2$
and one zero score.  Replacing $m$ by $4m$ scales $S(ma^2)$ by $1/4$ while
leaving $S(0)=1$ fixed, so the normalized weight on $a$, and hence the
output, changes, consistently with the one-layer identifiability theorem of
\citet{henry2025}.  In the depth-two construction, the first layer sends
every nonzero input to two nonzero tokens, so replacing $m_2$ by $4m_2$
scales every unnormalized second-layer weight by $1/4$ and preserves the
normalized weights.
\end{myremark}

\begin{myremark}
The normalizer in Theorem~\ref{thm:negative} satisfies all of the assumptions
on $S$ in Conjecture~3.10 of \citet[Section~3.4]{henry2025}: it is positive,
injective, and satisfies $S(0)=1$.  It is also Borel measurable.  The normalizer
is discontinuous at zero and unbounded near zero, so the counterexample does
not establish nonidentifiability for smooth normalizers; we leave open whether
continuity together with the hypotheses of \citet{henry2025} is sufficient
for generic identifiability.
\end{myremark}

\begin{theorem}[Identifiability for exponential softmax]
\label{thm:positive}
Let $d,d_{\mathrm{out}}\ge1$, $t\ge2$, $\ell\ge1$, and $\tau>0$, and set
$S_\tau(z)=e^{z/\tau}$.  If
$W=(M_1,\ldots,M_\ell,L)\in\cP_{d,d_{\mathrm{out}},\ell}$ satisfies a
nonvanishing condition on its final map and non-skewness conditions on the
score matrices before its final attention layer,
\begin{equation}
 L\ne0
 \quad\text{and}\quad M_r^\top\ne-M_r\quad\text{for }1\le r<\ell,
 \label{eq:source-regular}
\end{equation}
then
\begin{equation}
 \phi_{W'}^{S_\tau}=\phi_W^{S_\tau}
 \quad\Longrightarrow\quad W'=W,
 \label{eq:positive-result}
\end{equation}
for every $W'\in\cP_{d,d_{\mathrm{out}},\ell}$, without a nondegeneracy
assumption on $W'$.  Moreover, for any $1\le i\le d$ and
$1\le o\le d_{\mathrm{out}}$, the set
\begin{equation}
 \cU_{oi}
 =
 \left\{
 W\in\cP_{d,d_{\mathrm{out}},\ell}:
 L_{oi}\prod_{r=1}^{\ell-1}(M_r)_{ii}\ne0
 \right\},
 \label{eq:coordinate-locus}
\end{equation}
is nonempty, Euclidean open and dense, is contained in the set
specified by Equation~\eqref{eq:source-regular}, and satisfies
$\cF_{S_\tau}(W)=\{W\}$ for every $W\in\cU_{oi}$;
here the product is understood to equal one when $\ell=1$.
\end{theorem}

\begin{proof}
Without loss of generality, let $\tau=1$.  Suppose
$\phi_W^{S_\tau}=\phi_{W'}^{S_\tau}$, with
$W'=(M_1',\ldots,M_\ell',L')$ otherwise arbitrary.
We call the network with virtual parameters $W$ the source network and
the network with virtual parameters $W'$ the competing network.

Our identification of the virtual weights has five steps, with $M_r$ and
$M_r'$ compared successively for $r=1,\ldots,\ell$, namely,
1) constant token sequences determine the final linear map because every
attention layer fixes identical tokens and the network output repeats their
image under that map;
2) once $M_j=M_j'$ for every $j<r$,
Lemma~\ref{lem:two-class-closure} keeps the two token classes of
$(x+\varepsilon u,x-\varepsilon u,\ldots,x-\varepsilon u)$,
$\varepsilon\ne0$ and $u\ne0$, distinct through layers $1,\ldots,r-1$
whenever $u^\top M_j u\ne0$ for every $j<r$;
3) at layer $r$, Lemma~\ref{lem:quadratic-signal} computes the
quadratic coefficient of the displacement between corresponding output
tokens of the two networks and bounds each network's internal token gap at
cubic order in the incoming gap;
4) Lemma~\ref{lem:suffix-persistence} preserves this quadratic coefficient
through arbitrary, possibly different score matrices at layers
$r+1,\ldots,\ell$ because each
normalized layer takes convex combinations within its input token segment;
5) Lemma~\ref{lem:usable-directions} supplies a dense set of directions
$u$ satisfying $Lu\ne0$ and $u^\top M_j u\ne0$ for every $j<r$, so equality
of the realized functions determines $M_r$ successively for
$r=1,\ldots,\ell$;

Since constant token sequences pass unchanged through normalized attention layers 
\begin{equation}
 \phi_W^{S_\tau}(x,\ldots,x)
 =(Lx,\ldots,Lx)
 \quad\text{and}\quad
 \phi_{W'}^{S_\tau}(x,\ldots,x)
 =(L'x,\ldots,L'x).
 \label{eq:constant-recovery-proof}
\end{equation}
Thus, evaluating on repeated copies of each standard basis vector $e_j$ gives
$Le_j=L'e_j$ and establishes $L=L'$ by identifying every column.

Fix $1\le r\le\ell$ and suppose inductively that $M_j=M_j'$ for every
$j<r$.  To compare $M_r$ and $M_r'$, take an arbitrary base point
$x\in\R^d$ and choose, using Lemma~\ref{lem:usable-directions}, any
direction $u\in\R^d$ satisfying $Lu\ne0$ and $u^\top M_j u\ne0$ for
every $j<r$.  For a nonzero perturbation
$\varepsilon$, consider
\begin{equation}
 X_\varepsilon(x,u)
 =(x+\varepsilon u,x-\varepsilon u,\ldots,x-\varepsilon u),
 \label{eq:centered-probe-proof}
\end{equation}
with one token in the distinguished class and $t-1$ tokens in the repeated
class.  The condition $Lu\ne0$ prevents annihilation of nonzero
displacements along $u$, while Lemma~\ref{lem:two-class-closure} uses
$u^\top M_j u\ne0$ for $j<r$ to keep the two token classes distinct
through layers $1,\ldots,r-1$.

For token values $x+pu$ and $x+qu$,
Lemma~\ref{lem:two-class-closure} gives the layer action
\begin{equation}
 \T_M^{S_\tau}(x+pu,x+qu,\ldots,x+qu)
 =(x+p^+u,x+q^+u,\ldots,x+q^+u),
 \label{eq:two-class-form-proof}
\end{equation}
with the base point, direction, and multiplicities unchanged, and its
noncollapse conclusion applies to each $M_j$ with $j<r$ because
$u^\top M_j u\ne0$.  The common layers $1,\ldots,r-1$ therefore produce the same
scalar coefficients $p_\varepsilon$ and $q_\varepsilon$ in both networks
immediately before layer $r$, with $p_\varepsilon$ belonging to the
distinguished class and $q_\varepsilon$ to the repeated class.  Convexity
keeps both coefficients in $[-|\varepsilon|,|\varepsilon|]$, while prefix
noncollapse makes their difference $s_\varepsilon=p_\varepsilon-q_\varepsilon$
a nonzero incoming gap that tends to zero with $\varepsilon$.

In the source network, the layer-$r$ tokens have the form
$x+y_i u$, while those in the competing network have the form
$x+y_i' u$, with $i=1$ labeling the distinguished class and
$i=2$ the repeated class; both coefficient pairs depend on $\varepsilon$
through $p_\varepsilon$ and $q_\varepsilon$.  For each $i\in\{1,2\}$, the
quantitative estimates following Lemma~\ref{lem:quadratic-signal} express
the difference between corresponding scalar coefficients as
\begin{equation}
 y_i-y_i'
 =\kappa_n x^\top(M_r-M_r')u\,s_\varepsilon^2
  +O(|\varepsilon|s_\varepsilon^2),
 \label{eq:quadratic-signal-proof}
\end{equation}
where $\kappa_n=n/(1+n)^2$ is the positive multiplicity factor in the
quadratic term and $n=t-1$ is the repeated-class multiplicity.
Both classes have the same quadratic coefficient, although their
$O(|\varepsilon|s_\varepsilon^2)$ remainders can differ.  The difference
between the distinguished and repeated coefficients in the source network
satisfies
\begin{equation}
 y_1-y_2
 =\kappa_n(u^\top M_r u)s_\varepsilon^3+O(|s_\varepsilon|^4);
 \label{eq:cubic-gap-proof}
\end{equation}
likewise, the internal gap in the competing network satisfies
\begin{equation}
 y_1'-y_2'
 =\kappa_n(u^\top M_r' u)s_\varepsilon^3+O(|s_\varepsilon|^4),
 \label{eq:cubic-gap-primed-proof}
\end{equation}
so both internal gaps are $O(|s_\varepsilon|^3)$ and subdominant to the
quadratic scale in Equation~\eqref{eq:quadratic-signal-proof}.

In each network, every later token remains in the segment spanned by its
two layer-$r$ token values.  The $O(|s_\varepsilon|^3)$ bound on each
segment's diameter therefore also bounds the change in displacement between
corresponding tokens through arbitrary remaining layers.  Since
$p_\varepsilon$ and $q_\varepsilon$ lie in $[-|\varepsilon|,|\varepsilon|]$,
their gap satisfies $|s_\varepsilon|\le2|\varepsilon|$, which absorbs this
contribution into the $O(|\varepsilon|s_\varepsilon^2)$ remainder in the
quadratic expansion.  For the distinguished first output token, denoted by
the subscript $1$ and taking values in $\R^{d_{\mathrm{out}}}$, applying the
common final map therefore yields
\begin{equation}
 [\phi_W^{S_\tau}(X_\varepsilon)]_1
 -[\phi_{W'}^{S_\tau}(X_\varepsilon)]_1
 =s_\varepsilon^2
   \left[\kappa_n x^\top(M_r-M_r')u+O(|\varepsilon|)\right]Lu,
 \label{eq:directional-recovery-proof}
\end{equation}
whose left-hand side vanishes by equality of the realized functions.
Since $s_\varepsilon\ne0$ and $Lu\ne0$, a nonzero value of the
$\varepsilon$-independent coefficient $\kappa_n x^\top(M_r-M_r')u$
would dominate the $O(|\varepsilon|)$ remainder as $\varepsilon\to0$
and contradict functional equality; positivity of $\kappa_n$ therefore yields
\begin{equation}
 x^\top(M_r-M_r')u=0,
 \label{eq:directional-bilinear-proof}
\end{equation}
for every $x\in\R^d$ and every $u$ satisfying $Lu\ne0$ and
$u^\top M_j u\ne0$ for every $1\le j<r$.
For the polynomial
\begin{equation}
 P_r(u)=\|Lu\|^2\prod_{j=1}^{r-1}u^\top M_j u,
 \label{eq:usable-polynomial-proof}
\end{equation}
the set $\{u\in\R^d:P_r(u)\ne0\}$ consists exactly of the directions
satisfying $Lu\ne0$ and $u^\top M_j u\ne0$ for every $1\le j<r$.
Each factor of $P_r$ is not identically zero because $L\ne0$ and
$M_j^\top\ne-M_j$ for every $1\le j<r$.  By Lemma~\ref{lem:usable-directions},
$\{u:P_r(u)\ne0\}$ is dense, so for each fixed $x$, continuity in $u$ extends the
bilinear identity to every direction.  Taking $x=e_i$ and $u=e_j$ then
identifies each entry of $M_r-M_r'$, establishing $M_r=M_r'$ for the
inductive step.  At $r=1$, the prefix is empty, so the same argument starts
the induction without a preceding score condition; applying it through
layer $\ell$ establishes $W=W'$ by recovering every score matrix and the
final map.  Recovery of $M_r$ uses non-skewness only of $M_j$ for $j<r$,
so $M_\ell$ and all competing score matrices $M_1',\ldots,M_\ell'$
remain unrestricted.

\begingroup\emergencystretch=3em
It remains to prove the explicit open-dense clause.  The coordinate product
$L_{oi}\prod_{r=1}^{\ell-1}(M_r)_{ii}$ is a polynomial in the entries of
$W$ and is not identically zero, since it equals one when every matrix
entry of $W$ equals one.  Lemma~\ref{lem:usable-directions} therefore shows
that its nonvanishing set $\cU_{oi}$ is nonempty, Euclidean open, and dense.
If $W\in\cU_{oi}$, then $L_{oi}\ne0$
and $(M_r)_{ii}\ne0$ for every $r<\ell$, which ensures $L\ne0$ and excludes
skew-symmetry of each preceding score matrix because a real skew-symmetric
matrix has zero diagonal.  The pointwise identifiability result therefore
applies at every $W\in\cU_{oi}$, giving singleton complete fibers on this
open dense set.\par\endgroup

For any fixed primitive architecture whose query, key, and value matrices
have positive row and column dimensions,
substituting the query, key, and value matrices into the coordinate product
in Equation~\eqref{eq:coordinate-locus} gives a nonzero polynomial:
choosing every primitive matrix entry strictly positive makes $L_{oi}$ and
every $(M_r)_{ii}$ positive.  The resulting polynomial has an open dense
nonvanishing set in primitive parameter space, whose image under the
primitive-to-virtual map is precisely the intersection of $\cU_{oi}$ with
the architecture's virtual parameter image.  Continuity of the
primitive-to-virtual map makes this intersection dense in the virtual
parameter image, and openness of $\cU_{oi}$ makes it relatively open;
hence singleton complete fibers hold on a relatively open dense subset of
the virtual parameters realized by that architecture.
\end{proof}

\begin{myremark}
Although the directions used to recover $M_r$ are selected through the quadratic
forms $u^\top M_j u$ for $1\le j<r$, recovery uses the bilinear identity
$x^\top(M_r-M_r')u=0$ with $x$ and $u$ independent.  The argument therefore
determines the full score matrix, including its skew-symmetric component,
rather than only its symmetric part.
\end{myremark}

\begin{myremark}
Theorem~\ref{thm:positive} imposes one-sided hypotheses.  No nondegeneracy is
assumed of $W'$ a priori: constant inputs force $L'=L$, while recovery of
$M_r$ requires no nondegeneracy assumption on any unrecovered score matrix of
$W'$.  On the source side, non-skewness is required only for
$M_1,\ldots,M_{\ell-1}$ because recovery of $M_r$ requires noncollapse
only through the prefix $M_1,\ldots,M_{r-1}$; the terminal score $M_\ell$
belongs to none of these prefixes.  Likewise $L\ne0$ ensures that some displacement survives
the final map, without requiring $L$ to have full rank.

The full source set defined by Equation~\eqref{eq:source-regular} is itself
Euclidean open and dense: $L=0$ and each skew-symmetric matrix space are
proper closed linear subspaces of their respective finite-dimensional
ambient spaces.  $\cU_{oi}$ is a more explicit algebraic open-dense
subset obtained by choosing one coordinate direction that satisfies every
required nonvanishing condition simultaneously.  
\end{myremark}

\section{Tool and computational resource disclosure}
\label{sec:ai-disclosure}

Following the disclosure and verifiability principles advocated by the
Leiden Declaration on Artificial Intelligence and Mathematics
\citep{leiden2026} and discussed by \citet{tao2026ai}, we record the role of
automated tools in this work. Codex was used with Lean~4 for proof search,
theorem proving, proof debugging, inspection and summarization of formal code,
and manuscript drafting and revision. The authors report extensive use of
GPT-5.6 Sol, primarily at very-high reasoning effort, with aggregate inference
on the order of tens of millions of tokens. This is an author-reported estimate,
not a measurement of the verification commands.

James Amarel is the formalizer and maintainer. The mathematical claims and
exposition remain the responsibility of the authors. Supporting Lean code,
including the headline theorems, log-affine corollaries, and boundary examples,
is organized in the DeepLean root Lake project.\footnote{Code: \url{https://github.com/lanl/DeepLean}.} Its Challenge/Solution
interface exposes the three declarations described in
Appendix~\ref{app:formal-certification}. Local compilation, statement comparison,
transitive axiom audits, and independent proof replay have distinct roles;
none establishes the correspondence between formal statements and mathematical
prose without a separate semantic review. The reproducible commands are
documented in the repository's \path{README.md}; each run records its results
under \path{.cache/local-proof/}.

\section{Conclusion}
\label{sec:conclusion}

Positivity and injectivity of the scalar normalizer do not suffice for
identifiability of deep normalized attention.
Indeed, row normalization retains ratios rather
than individual weights, and we exhibit an admissible normalizer whose ratio
structure produces non-singleton virtual fibers on an open dense set.
Exponential softmax behaves differently: at every finite depth, complete
virtual fibers are singletons under the explicit conditions of
Theorem~\ref{thm:positive}.  Injectivity of $S$ permits recovery of a scalar
score $z$ from $S(z)$, but does not by itself exclude changes in the score
matrices that, for every input, multiply the unnormalized weights within
each attention row by a common positive factor and hence preserve the
realized function.

For softmax, Theorem~\ref{thm:positive} replaces the local evidence from
finite-batch Jacobian ranks for deep normalized attention with a global statement on the complete functional
fiber.  Equality of the realized function on the whole input space recovers
the virtual score matrices successively, even when the competing parameter
point is arbitrary within the same virtual parameter space, and the hypotheses
hold on an explicit nonempty Euclidean-open dense set.  Thus normalization cannot be understood simply as
removing a rescaling symmetry from lightning attention: its ratio structure
can either create persistent functional ambiguity or eliminate it under
composition.  In the softmax case considered here, the latter rigidity
separates the remaining raw query-key and value-coordinate redundancies from
the virtual identifiability of the deep normalized model.
\clearpage
\bibliography{pmlr-sample}

\clearpage
\appendix

\section{The counterexample normalizer and its fiber}
\label{app:counterexample}

This appendix verifies positivity, injectivity, and Borel measurability of
the normalizer in Theorem~\ref{thm:negative} and proves that
$(m_1,m_2,l)$ and $(m_1,4m_2,l)$ realize the same function whenever
$m_1m_2l\ne0$.

\subsection{Injectivity, measurability, and discontinuity}

Write $\{u\}=u-\lfloor u\rfloor$ for the fractional part of a real number
$u$.  For $x>0$, Equation~\eqref{eq:exotic-normalizer} gives
$\{g(x)\}=3/4+r(x)/8\in[3/4,7/8)$, whereas for $x<0$ the decomposition
\begin{equation}
 g(x)=(-k(x)-1)+\left(\frac14+\frac18r(x)\right),
 \label{eq:negative-g-decomposition}
\end{equation}
places $\{g(x)\}$ in $[1/4,3/8)$.  The disjoint intervals determine the
sign of $x$ from $\{g(x)\}$, while their exclusion of zero separates
$g(0)=0$ from every $g(x)$ with $x\ne0$.

For $x>0$, the integer part of $g(x)$ recovers
$k(x)=-\lfloor g(x)\rfloor$, and its fractional part recovers
$r(x)=8(\{g(x)\}-3/4)$.  For $x<0$, the integer shift in
Equation~\eqref{eq:negative-g-decomposition} gives
$k(x)=-\lfloor g(x)\rfloor-1$, while the fractional part yields
$r(x)=8(\{g(x)\}-1/4)$.  For $x\ne0$, reconstruction of the magnitude follows from
\begin{equation}
 |x|=4^{k(x)+r(x)},
 \label{eq:g-decoding}
\end{equation}
which determines $x$ together with its recovered sign and proves
injectivity of $g$.  Since $u\mapsto4^u$ is strictly increasing, $S=4^g$
is injective, and positivity and $S(0)=1$ follow from its definition.

On each half-line $x>0$ and $x<0$, the function $g$ is obtained from
continuous functions, the Borel floor function, and arithmetic operations.
Assigning $g(0)=0$ on the remaining singleton therefore makes $g$ Borel
measurable on $\R$, and continuity of exponentiation makes $S=4^g$ Borel
measurable.  The sequence $x_n=4^{-n}$ tends to zero and satisfies
$k(x_n)=-n$ with $r(x_n)=0$, so $S(x_n)=4^{n+3/4}\to\infty$ while
$S(0)=1$, proving discontinuity at zero.

Positivity, injectivity, and $S(0)=1$ are the normalizer hypotheses of
\citet[Section~3.4]{henry2025}; Borel measurability is an additional
property of the construction.

\subsection{Scaling and domination}

For $x\ne0$, the identity $q(4x)=q(x)+1$ gives $k(4x)=k(x)+1$ and
leaves $r(x)$ unchanged.  Since multiplication by four preserves the sign
of $x$, the definition of $g$ yields $g(4x)=g(x)-1$, and exponentiation
gives
\begin{equation}
 S(4x)=\frac14S(x),
 \label{eq:four-scaling-app}
\end{equation}
so scaling all scores in a row by four leaves its normalized weights
unchanged whenever every score in that row is nonzero.  At $z=0$, the
equation $S(4z)=S(z)/4$ fails because $S(0)=1$, so using a uniform
factor $1/4$ for all second-layer weights requires every second-layer
score to be nonzero.

For arbitrary $z,w>0$, the construction also satisfies
\begin{equation}
 zS(z)>wS(-w),
 \label{eq:mass-dominance}
\end{equation}
because the positive and negative branches of $g$ give separated
exponent ranges.  Indeed, multiplying $S(z)$ by $z=4^{k(z)+r(z)}$ yields
\begin{equation}
 zS(z)=4^{3/4+9r(z)/8},
 \label{eq:mass-identities}
\end{equation}
whose exponent is at least $3/4$, whereas multiplying $S(-w)$ by
$w=4^{k(w)+r(w)}$ yields $wS(-w)=4^{-3/4+9r(w)/8}$, whose exponent is
strictly less than $3/8$.

\subsection{Nonvanishing of the first-layer outputs}

Fix a nonzero score parameter $m$ and a nonzero input $(a,b)$, and write
$(y_1,y_2)=\T_m^S(a,b)$.  The first output is
\begin{equation}
 y_1=
 \frac{S(ma^2)a+S(mab)b}
      {S(ma^2)+S(mab)},
 \label{eq:scalar-outputs}
\end{equation}
whose denominator is positive because $S$ is positive; exchanging $a$
and $b$ exchanges the two outputs, so nonvanishing of the numerator of
$y_1$ for every nonzero input also proves $y_2\ne0$.  If $ab\ge0$, the
two terms $aS(ma^2)$ and $bS(mab)$ are both nonnegative or both
nonpositive, and at least one is nonzero, so their sum cannot vanish.

For $ab<0$, simultaneous negation of $a$ and $b$ preserves the scores
$ma^2$ and $mab$ and negates the numerator of $y_1$, so it suffices to
consider $a>0>b$.  When $m>0$, applying
Equation~\eqref{eq:mass-dominance} to $z=ma^2>0$ and $w=-mab>0$ and
canceling the positive factor $ma$ gives
\begin{equation}
 aS(ma^2)>(-b)S(mab),
 \label{eq:positive-m-dominance}
\end{equation}
so the numerator $aS(ma^2)+bS(mab)$ of $y_1$ is positive.  When $m<0$, the positive
score is $mab$ and the negative score is $ma^2$; applying
Equation~\eqref{eq:mass-dominance} with $z=mab$ and $w=-ma^2$ and
canceling the positive factor $-ma$ gives
\begin{equation}
 (-b)S(mab)>aS(ma^2),
 \label{eq:negative-m-dominance}
\end{equation}
so the numerator $aS(ma^2)+bS(mab)$ of $y_1$ is negative.

For every $m\ne0$ and every nonzero input $(a,b)$, the sign analysis and
the exchange of the two input tokens establish
\begin{equation}
 y_1y_2\ne0,
 \label{eq:first-layer-nonvanishing}
\end{equation}
so every second-layer score $m_2y_i y_j$ is nonzero for
$i,j\in\{1,2\}$ when $m=m_1$ and $m_1m_2\ne0$.

\subsection{The non-singleton fiber}

Let $W=(m_1,m_2,l)$ with $m_1m_2l\ne0$, and set
$W^\sharp=(m_1,4m_2,l)$.  Since both networks have first-layer score
parameter $m_1$, they produce the same first-layer output $(y_1,y_2)$
on each input $(a,b)$.  For every nonzero input,
Equation~\eqref{eq:first-layer-nonvanishing} and $m_2\ne0$ make each
second-layer score $m_2y_i y_j$ nonzero, so
Equation~\eqref{eq:four-scaling-app} gives
\begin{equation}
 S(4m_2y_i y_j)=\frac14S(m_2y_i y_j)
 \quad\text{for }i,j\in\{1,2\},
 \label{eq:second-layer-scaling}
\end{equation}
so each row's unnormalized weights acquire the common factor $1/4$,
which cancels under normalization and leaves both second-layer outputs
unchanged.  At $(a,b)=(0,0)$, both first-layer tokens and both second-layer
tokens equal zero, so both realizations vanish.  Applying the common
final map $l$ therefore yields
\begin{equation}
 \phi_{W^\sharp}^S=\phi_W^S,
 \label{eq:counterexample-fiber-equality}
\end{equation}
with $W^\sharp\ne W$ because $m_2\ne0$, proving that $\cF_S(W)$ is
non-singleton.

The set $\{m_1m_2l\ne0\}\subset\R^3$ is the complement of the three
coordinate hyperplanes, and is therefore nonempty, Euclidean open, and
dense.

\section{Supporting results for Theorem~\ref{thm:positive}}
\label{app:asymptotics}

Constant inputs determine the final map through
Equation~\eqref{eq:constant-recovery-proof}; recovery of each score matrix
then requires a nonzero perturbation that survives the preceding layers and
produces a displacement that the remaining layers cannot cancel.  We retain
$S_\tau(z)=e^{z/\tau}$ with $\tau=1$, as fixed in the proof of
Theorem~\ref{thm:positive}.  The four lemmas below establish noncollapse of
the two token classes, the quadratic displacement between networks and cubic
gap within each network, persistence through arbitrary remaining score
matrices, and the density of directions on which the displacement survives
the final map.

For a query with one distinguished key and $n\ge1$ identical repeated keys,
the normalized weight on the distinguished key depends on its score
difference $z$ from a repeated key through
\begin{equation}
 \sigma_n(z)=\frac{e^z}{e^z+n},
 \label{eq:weighted-sigmoid}
\end{equation}
which lies strictly between zero and one because every unnormalized weight
is positive.  The factor $n$ records the repeated keys' total contribution
to the normalizing denominator.

\subsection{Two-class closure and noncollapse}

\begin{mylemma}[Two-class closure and noncollapse]
\label{lem:two-class-closure}
Let $t\ge2$ and $n=t-1$, and let $M\in\mathbb R^{d\times d}$,
$x,u\in\mathbb R^d$ with $u\ne0$, and $p,q\in\mathbb R$.  The input
\begin{equation}
 X(x,u;p,q)=(x+pu,x+qu,\ldots,x+qu),
 \label{eq:two-class-input}
\end{equation}
has one distinguished token and $n$ repeated tokens.  Write $s=p-q$ for
the signed difference of their scalar coefficients, and set $a=x^\top Mu$
and $b=u^\top Mu$, so the distinguished query's score difference is
$s(a+pb)$.  Its output coefficient is
\begin{equation}
 p^+=q+\sigma_n\!\bigl(s(a+pb)\bigr)s,
 \label{eq:two-class-p-update}
\end{equation}
where the sigmoid factor is the normalized weight on the distinguished
input.  Every repeated query instead has output coefficient
\begin{equation}
 q^+=q+\sigma_n\!\bigl(s(a+qb)\bigr)s,
 \label{eq:two-class-q-update}
\end{equation}
because its query token is $x+qu$.  Consequently the layer action is
\begin{equation}
 \T_M^{S_\tau}X(x,u;p,q)=X(x,u;p^+,q^+),
 \label{eq:two-class-exact-closure}
\end{equation}
with the base point, direction, and multiplicities preserved, and both
output coefficients between $p$ and $q$.  If $s\ne0$ and $b\ne0$, then
$p^+\ne q^+$, so the two output token values remain distinct.

For any finite list $M_1,\ldots,M_k$ with $k\ge0$ and
$u^\top M_j u\ne0$ for every $1\le j\le k$, the centered input
$X(x,u;\varepsilon,-\varepsilon)$ therefore retains distinct token values
through all $k$ layers whenever $\varepsilon\ne0$.  Both final scalar
coefficients $p_\varepsilon$ and $q_\varepsilon$ lie in
$[-|\varepsilon|,|\varepsilon|]$, so they tend to zero as
$\varepsilon\to0$ while their difference remains nonzero for every
$\varepsilon\ne0$.
\end{mylemma}

\begin{proof}
For query $x+pu$, subtracting the score of a repeated key from the score of
the distinguished key gives
\begin{equation}
 (x+pu)^\top M(x+pu)-(x+pu)^\top M(x+qu)=s(a+pb),
 \label{eq:two-class-distinguished-score}
\end{equation}
which determines the ratio between their unnormalized weights.  Factoring
$e^{(x+pu)^\top M(x+qu)}$ from the attention row gives distinguished-key
weight $\sigma_n(s(a+pb))$; averaging the two token values with this weight
gives Equation~\eqref{eq:two-class-p-update}.  For query $x+qu$, the score
difference is $s(a+qb)$, so the same normalization gives
Equation~\eqref{eq:two-class-q-update}.  All repeated queries are identical
and hence produce the same output token, establishing
Equation~\eqref{eq:two-class-exact-closure}.

Subtracting the two coefficient updates gives
\begin{equation}
 p^+-q^+
 =s\left[\sigma_n\!\bigl(s(a+pb)\bigr)
          -\sigma_n\!\bigl(s(a+qb)\bigr)\right],
 \label{eq:two-class-gap-update}
\end{equation}
so noncollapse follows from distinct sigmoid arguments whenever $s\ne0$.
The arguments differ by
\begin{equation}
 s(a+pb)-s(a+qb)=s^2b,
 \label{eq:two-class-logit-gap}
\end{equation}
which is nonzero when $s\ne0$ and $b\ne0$.  For $n>0$, equality
$\sigma_n(z)=\sigma_n(z')$ implies $ne^z=ne^{z'}$ by cross multiplication,
and hence $z=z'$.  Injectivity of $\sigma_n$ therefore makes the bracket in
Equation~\eqref{eq:two-class-gap-update} nonzero, proving $p^+\ne q^+$.

For the finite list, let $p_j(\varepsilon)$ and $q_j(\varepsilon)$ denote
the distinguished and repeated coefficients after $j$ layers, starting
with $p_0(\varepsilon)=\varepsilon$ and $q_0(\varepsilon)=-\varepsilon$.
Applying the exact noncollapse result
successively with $b=u^\top M_j u$ gives
\begin{equation}
 p_j(\varepsilon)-q_j(\varepsilon)\ne0
 \quad\text{for }0\le j\le k\text{ and }\varepsilon\ne0,
 \label{eq:two-class-prefix-noncollapse}
\end{equation}
so the two token values remain distinct at every layer.  Each coefficient update is a
convex combination of the incoming coefficients because
$0<\sigma_n(z)<1$.  The coefficient intervals are therefore nested inside
$[-|\varepsilon|,|\varepsilon|]$, proving convergence of both coefficients
to zero after every layer, including the empty list $k=0$.
\end{proof}

At step $r$ of Theorem~\ref{thm:positive}, equality $M_j=M_j'$ for every
$j<r$ gives both networks the same coefficients $p_\varepsilon$ and
$q_\varepsilon$ after applying the common prefix to
$X(x,u;\varepsilon,-\varepsilon)$.  For $\varepsilon\ne0$ and a direction satisfying $Lu\ne0$ and
$u^\top M_j u\ne0$ for every $j<r$, Lemma~\ref{lem:two-class-closure}
ensures that their signed difference $s_\varepsilon=p_\varepsilon-q_\varepsilon$
is nonzero and satisfies $|s_\varepsilon|\le2|\varepsilon|$, so it tends
to zero while the token values remain distinct.  The parameter
$\varepsilon$ specifies the original input perturbation, whereas
$s_\varepsilon$ measures the gap on which the comparison of $M_r$ and
$M_r'$ depends.

\subsection{Quadratic signal and cubic gap}

\begin{mylemma}[Quadratic signal and cubic gap]
\label{lem:quadratic-signal}
Let $t\ge2$ and $n=t-1$, and fix $x,u\in\mathbb R^d$ with $u\ne0$ and
$N,N'\in\mathbb R^{d\times d}$.  Suppose $p_\varepsilon$ and
$q_\varepsilon$ tend to zero, with signed gap
$s_\varepsilon=p_\varepsilon-q_\varepsilon\ne0$ for every sufficiently
small nonzero $\varepsilon$.  Apply $\T_N^{S_\tau}$ to
$X(x,u;p_\varepsilon,q_\varepsilon)$ and denote its distinguished output by
$x+y_1u$ and its repeated output by $x+y_2u$; use $y_1'$ and $y_2'$ for
the corresponding coefficients after applying $\T_{N'}^{S_\tau}$ to the
same input.  These coefficients depend on $\varepsilon$ and are uniquely
determined because $u\ne0$.

Write $a=x^\top Nu$ and $a'=x^\top N'u$ for the two matrices' scores
between the base point and direction, and write $b=u^\top Nu$ and
$b'=u^\top N'u$ for their quadratic scores on the direction.  All four
output coefficients tend to zero.  The positive coefficient
$\kappa_n=\sigma_n'(0)=n/(1+n)^2$ determines the displacement for each
corresponding token class through
\begin{equation}
 y_i-y_i'=\kappa_n(a-a')s_\varepsilon^2+o(s_\varepsilon^2)
 \quad\text{for }i\in\{1,2\},
 \label{eq:quadratic-signal-limit}
\end{equation}
where both classes have the same quadratic coefficient, although their
remainders can differ.  Within the output of the layer with matrix $N$, the
class gap satisfies
\begin{equation}
 y_1-y_2=\kappa_n b\,s_\varepsilon^3+O(|s_\varepsilon|^4),
 \label{eq:quantitative-cubic-gap}
\end{equation}
whereas the layer with matrix $N'$ satisfies
\begin{equation}
 y_1'-y_2'=\kappa_n b'\,s_\varepsilon^3+O(|s_\varepsilon|^4),
 \label{eq:cubic-gap-primed-limit-proof}
\end{equation}
so both internal gaps are of smaller order than the cross-network
quadratic term whenever $a\ne a'$, including when $b$ or $b'$ vanishes.
\end{mylemma}

\begin{proof}
For $i\in\{1,2\}$, let $c_1=p_\varepsilon$ denote the distinguished input
coefficient and $c_2=q_\varepsilon$ the repeated input coefficient.
Lemma~\ref{lem:two-class-closure} gives
\begin{equation}
 y_i=q_\varepsilon+
       \sigma_n\!\bigl(s_\varepsilon(a+c_i b)\bigr)s_\varepsilon,
 \label{eq:quadratic-signal-unprimed-update}
\end{equation}
because the sigmoid factor is the normalized weight on the distinguished
key for that class's query.  Replacing $(a,b)$ by $(a',b')$ gives the
coefficient $y_i'$ for the layer with matrix $N'$.  Each output coefficient
lies between $p_\varepsilon$ and $q_\varepsilon$, so all four tend to zero.

The weighted sigmoid has derivative
\begin{equation}
 \sigma_n'(z)=\frac{ne^z}{(e^z+n)^2},
 \label{eq:quadratic-signal-sigmoid-derivative}
\end{equation}
whose value at zero is $\kappa_n$.  Since $\sigma_n''$ is bounded on a
neighborhood of zero, integrating $\sigma_n'$ between arguments $A$ and
$B$ in that neighborhood gives
\begin{equation}
 \sigma_n(A)-\sigma_n(B)
 =(A-B)\left[\kappa_n+O(|A|+|B|)\right],
 \label{eq:quantitative-sigmoid-difference}
\end{equation}
with the same identity when $A=B$.  All arguments in the coefficient
updates are $O(|s_\varepsilon|)$ because $p_\varepsilon$ and
$q_\varepsilon$ tend to zero.  For corresponding classes, the argument
difference is $s_\varepsilon[(a-a')+c_i(b-b')]$, so the outer factor
$s_\varepsilon$ in Equation~\eqref{eq:quadratic-signal-unprimed-update}
gives
\begin{equation}
 y_i-y_i'
 =\kappa_n[(a-a')+c_i(b-b')]s_\varepsilon^2
   +O(|s_\varepsilon|^3),
 \label{eq:quantitative-class-displacement}
\end{equation}
whose class-dependent quadratic correction is $o(s_\varepsilon^2)$ since
$c_i\to0$.  The cubic remainder is also $o(s_\varepsilon^2)$, proving
Equation~\eqref{eq:quadratic-signal-limit}.

For the two classes within the layer with matrix $N$, the sigmoid
arguments instead differ by
\begin{equation}
 s_\varepsilon(a+p_\varepsilon b)
 -s_\varepsilon(a+q_\varepsilon b)=s_\varepsilon^2b,
 \label{eq:cubic-gap-argument-difference}
\end{equation}
so Equation~\eqref{eq:quantitative-sigmoid-difference} gives their sigmoid
difference as $s_\varepsilon^2b[\kappa_n+O(|s_\varepsilon|)]$.
Multiplication by the outer $s_\varepsilon$ in the exact gap update
Equation~\eqref{eq:two-class-gap-update} proves
Equation~\eqref{eq:quantitative-cubic-gap}, with its signed cubic term and
fourth-order remainder.  Applying the same argument to $N'$ proves
Equation~\eqref{eq:cubic-gap-primed-limit-proof}; if either quadratic score
$b$ or $b'$ is zero, the corresponding exact internal gap is zero as well.
\end{proof}

For the centered probes in Theorem~\ref{thm:positive},
Lemma~\ref{lem:two-class-closure} places both incoming coefficients in
$[-|\varepsilon|,|\varepsilon|]$, so $c_i=O(|\varepsilon|)$ and
$|s_\varepsilon|\le2|\varepsilon|$.  The class-dependent correction and
cubic remainder in Equation~\eqref{eq:quantitative-class-displacement}
therefore give the sharper comparison
\begin{equation}
 y_i-y_i'
 =\kappa_n x^\top(N-N')u\,s_\varepsilon^2
  +O(|\varepsilon|s_\varepsilon^2),
 \label{eq:quantitative-quadratic-signal}
\end{equation}
whose remainder is dominated by the quadratic term whenever
$x^\top(N-N')u\ne0$.  Both internal class gaps remain
$O(|s_\varepsilon|^3)$, so their effect on later outputs is smaller than
this quadratic displacement.

\subsection{Persistence through unequal suffixes}

\begin{mylemma}[Persistence through unequal suffixes]
\label{lem:suffix-persistence}
Let $t\ge2$ and $x,u\in\mathbb R^d$ with $u\ne0$.  Suppose the scalar
coefficients $p_\varepsilon,q_\varepsilon,p'_\varepsilon,q'_\varepsilon$
all tend to zero, and let $\delta_\varepsilon\to0$ be nonzero for every
sufficiently small nonzero $\varepsilon$.  For a fixed $w\in\mathbb R$,
suppose the distinguished coefficients satisfy
$p_\varepsilon-p'_\varepsilon=w\delta_\varepsilon+o(|\delta_\varepsilon|)$
and the repeated coefficients satisfy
$q_\varepsilon-q'_\varepsilon=w\delta_\varepsilon+o(|\delta_\varepsilon|)$.
Assume also that $p_\varepsilon-q_\varepsilon=o(|\delta_\varepsilon|)$
within the unprimed network and
$p'_\varepsilon-q'_\varepsilon=o(|\delta_\varepsilon|)$ within the primed
network.

Apply layers with arbitrary score matrices $M,M'\in\mathbb R^{d\times d}$ to
$X(x,u;p_\varepsilon,q_\varepsilon)$ and
$X(x,u;p'_\varepsilon,q'_\varepsilon)$, respectively.  Denote the
distinguished and repeated output coefficients by $y_1$ and $y_2$ after
$\T_M^{S_\tau}$, and by $y_1'$ and $y_2'$ after $\T_{M'}^{S_\tau}$.
All four output coefficients tend to zero, and corresponding classes
satisfy
\begin{equation}
 y_i-y_i'=w\delta_\varepsilon+o(|\delta_\varepsilon|)
 \quad\text{for }i\in\{1,2\},
 \label{eq:suffix-persistence-output-corresponding}
\end{equation}
with internal gap $y_1-y_2=o(|\delta_\varepsilon|)$ in the unprimed
network and $y_1'-y_2'=o(|\delta_\varepsilon|)$ in the primed network.
The same conclusions hold after any two finite lists of softmax layers
of equal length, with independent score matrices in the two lists.
\end{mylemma}

\begin{proof}
Every normalized layer takes convex combinations of its input token
values.  By Lemma~\ref{lem:two-class-closure}, each network remains in the
two-class family, and its coefficient intervals are nested under further
layers.  Consequently, a scalar coefficient $z$ at any output position
after any number of unprimed layers satisfies
\begin{equation}
 |z-q_\varepsilon|\le|p_\varepsilon-q_\varepsilon|,
 \label{eq:quantitative-suffix-bounds}
\end{equation}
because its token stays between $x+p_\varepsilon u$ and
$x+q_\varepsilon u$.  At the corresponding position in the primed
network, the coefficient $z'$ satisfies
\begin{equation}
 |z'-q'_\varepsilon|\le|p'_\varepsilon-q'_\varepsilon|,
 \label{eq:quantitative-suffix-primed-bound}
\end{equation}
by containment in that network's own token segment, independently of the
score matrices used in either list.  Both bounds tend to zero, and their
reference coefficients $q_\varepsilon$ and $q'_\varepsilon$ also tend to
zero, proving convergence of every output coefficient.

The triangle inequality gives
\begin{equation}
 \bigl|(z-z')-(q_\varepsilon-q'_\varepsilon)\bigr|
 \le |p_\varepsilon-q_\varepsilon|
     +|p'_\varepsilon-q'_\varepsilon|,
 \label{eq:suffix-persistence-corresponding-proof}
\end{equation}
whose right-hand side is $o(|\delta_\varepsilon|)$ by the two internal-gap
hypotheses.  Since the repeated coefficients differ by
$w\delta_\varepsilon+o(|\delta_\varepsilon|)$, every corresponding pair
of output coefficients retains the coefficient $w$ at scale
$\delta_\varepsilon$, proving
Equation~\eqref{eq:suffix-persistence-output-corresponding}.  Any two
unprimed output coefficients differ by at most
$|p_\varepsilon-q_\varepsilon|$, while any two primed output coefficients
differ by at most $|p'_\varepsilon-q'_\varepsilon|$, so both internal gaps
remain $o(|\delta_\varepsilon|)$.  These interval bounds apply after
every finite number of layers and establish the asserted persistence
without a relation between the two lists' score matrices.
\end{proof}

For the layers with matrices $N$ and $N'$ in
Lemma~\ref{lem:quadratic-signal}, retain $y_1,y_2$ and $y_1',y_2'$ for
their output coefficients, and let $z,z'$ be coefficients at corresponding
token positions after the remaining layers.  Applying the interval bounds
to the two internal gaps, each $O(|s_\varepsilon|^3)$, gives
\begin{equation}
 z-z'=y_2-y_2'+O(|s_\varepsilon|^3),
 \label{eq:quantitative-suffix-persistence}
\end{equation}
because the remaining normalized layers keep the source tokens between
$x+y_1u$ and $x+y_2u$, and keep the competing tokens between $x+y_1'u$ and
$x+y_2'u$.  For centered probes,
$|s_\varepsilon|\le2|\varepsilon|$ absorbs this change into the
$O(|\varepsilon|s_\varepsilon^2)$ remainder in
Equation~\eqref{eq:quantitative-quadratic-signal}, preserving its quadratic
coefficient through the remaining layers.  A nonzero quadratic coefficient
$\kappa_n x^\top(N-N')u$ therefore cannot be canceled by either network's
smaller internal gap, even when all
remaining source and competing score matrices differ.


\subsection{Polynomial nonvanishing}

\begin{mylemma}[Polynomial nonvanishing]
\label{lem:usable-directions}
Let $m,d,d_{\mathrm{out}}\ge1$ and $k\ge0$ be integers.  For every
nonzero real polynomial $P$ on $\mathbb R^m$, the set
\begin{equation}
 \{z\in\mathbb R^m:P(z)\ne0\}
 \label{eq:polynomial-nonvanishing-set}
\end{equation}
is nonempty, Euclidean open, and dense.  In particular, if
$L\in\mathbb R^{d_{\mathrm{out}}\times d}$ is nonzero and
$M_1,\ldots,M_k\in\mathbb R^{d\times d}$ satisfy
$M_j^\top\ne-M_j$ for every $1\le j\le k$, then
\begin{equation}
 P(u)=\|Lu\|^2\prod_{j=1}^k u^\top M_j u
 \label{eq:usable-directions-polynomial}
\end{equation}
is a nonzero polynomial whose nonvanishing set consists precisely of the
directions satisfying $Lu\ne0$ and $u^\top M_j u\ne0$ for every
$1\le j\le k$.  For $k=0$, the product is one, so the conclusion concerns
only the condition $Lu\ne0$.
\end{mylemma}
\begin{proof}
Continuity of $P$ makes its nonvanishing set open; density follows once its
zero set is shown to have empty interior.  We prove the latter assertion by
induction on $m$.  A nonzero univariate polynomial has finitely many roots
and therefore cannot vanish on an open interval.  For $m>1$, write
$z=(v,s)\in\mathbb R^{m-1}\times\mathbb R$ and expand in the final
coordinate as
\begin{equation}
 P(v,s)=\sum_{j=0}^r a_j(v)s^j,
 \label{eq:polynomial-coordinate-expansion}
\end{equation}
where the coefficient polynomials $a_j$ include at least one nonzero
polynomial because $P$ is nonzero.  If $P$ vanished on a nonempty open
set, it would vanish on a product $U\times I$ with $U$ nonempty and open
in $\mathbb R^{m-1}$ and $I$ a nonempty open interval.  For each $v\in U$,
the polynomial $s\mapsto P(v,s)$ would then vanish identically, forcing
$a_j(v)=0$ for every $j$.  A nonzero coefficient polynomial would therefore
vanish on $U$, contradicting the induction hypothesis.  Thus the zero set
has empty interior, and its complement is dense and nonempty.

For the specialization in Equation~\eqref{eq:usable-directions-polynomial},
the factor $\|Lu\|^2$ is a nonzero polynomial because $L\ne0$.  Each
quadratic factor satisfies
\begin{equation}
 u^\top M_j u
 =u^\top\left(\frac{M_j+M_j^\top}{2}\right)u,
 \label{eq:quadratic-symmetric-part}
\end{equation}
so it depends only on the symmetric part of $M_j$.  If this quadratic
form vanished for every $u$, evaluation at each coordinate vector would
make every diagonal entry of the symmetric part zero, while evaluation at
sums of two coordinate vectors would make every off-diagonal entry zero.
Consequently $M_j+M_j^\top=0$, contrary to the hypothesis on $M_j$.
All factors of $P$ are therefore nonzero polynomials, and the absence of
zero divisors in the real polynomial ring makes their product nonzero.
Since $\|Lu\|^2$ vanishes exactly when $Lu=0$, nonvanishing of $P$ is
equivalent to simultaneous nonvanishing of $Lu$ and every $u^\top M_j u$.
Applying the first part of the lemma proves the asserted openness and
density, including $k=0$, when only $\|Lu\|^2$ remains.
\end{proof}

For recovery of $M_r$ in Theorem~\ref{thm:positive}, apply
Lemma~\ref{lem:usable-directions} to $L$ and the prefix
$M_1,\ldots,M_{r-1}$.  The condition $Lu\ne0$ preserves nonzero output
displacements along $u$, while $u^\top M_j u\ne0$ for $j<r$ prevents
collapse through the common prefix.  On this dense set of directions,
Equation~\eqref{eq:directional-bilinear-proof} gives
$x^\top(M_r-M_r')u=0$ for every $x\in\R^d$.
For each fixed $x$, continuity in $u$ extends the identity to every
direction, and independent variation of $x$ and $u$ determines every
entry of $M_r-M_r'$, including its skew-symmetric component.

\subsection{Boundary cases}
\label{app:boundaries}

If $t=1$, row normalization gives a single weight equal to one, so each
attention layer fixes its input token independently of the score matrix.
If $L=0$, every score list realizes the zero function.  Thus $t\ge2$ and
$L\ne0$ are necessary for the complete score recovery asserted in
Theorem~\ref{thm:positive}.

At two tokens, a skew-symmetric score matrix produces an additional
obstruction.  For $M^\top=-M$ and input $(x,y)$, write $a=x^\top My$;
the self-scores vanish and the reversed cross-score is $-a$.
Both output tokens coincide because the second satisfies
\begin{equation}
 \frac{e^{-a}x+y}{e^{-a}+1}
 =\frac{x+e^a y}{1+e^a},
 \label{eq:skew-two-token-collapse}
\end{equation}
whose right-hand side is the first output token, so the output sequence is
constant.
Every subsequent normalized layer fixes this constant sequence, regardless
of its score matrix.  In particular, taking
\begin{equation}
 M_1=\begin{pmatrix}0&1\\-1&0\end{pmatrix}
 \label{eq:full-rank-skew-example}
\end{equation}
and the final map $L=I_2$ gives distinct depth-two parameters
$W=(M_1,I_2,I_2)$ and $W'=(M_1,2I_2,I_2)$ with equal complete
realizations.  All score matrices in this example have full rank, so full
rank does not prevent the collapse.

The proper-prefix non-skew condition in Theorem~\ref{thm:positive}
ensures noncollapse on a dense set of probe directions; the two-token
example does not establish a converse classification of identifiable
parameters.  Recovery of $M_r$ requires noncollapse only through
$M_1,\ldots,M_{r-1}$, so the terminal matrix $M_\ell$ is unrestricted.
Likewise, $L\ne0$ supplies a dense set of directions satisfying $Lu\ne0$
without requiring full rank.  No nondegeneracy condition is imposed on
unrecovered competing score matrices, because the suffix bounds in
Lemma~\ref{lem:suffix-persistence} use only convexity.


\section{Formal verification and computational provenance}
\label{app:formal-certification}

\subsection{AI-assisted development}

Codex assists Lean proof search, theorem proving, debugging, inspection of
formal code, repository organization, and manuscript drafting. The disclosure
in Section~\ref{sec:ai-disclosure} records the human and tool contributions.
James Amarel is the formalizer and maintainer. 
\subsection{Statements and correspondence}

DeepLean uses the repository's root Lake project. Supporting
mathematics is in \path{DeepLean/Attention/}, with public imports in
\path{DeepLean.lean}.  Its principal version pins are
\begin{lstlisting}
Lean 4.35.0-rc2
Mathlib 4.35.0-rc2
  commit 065356127b1dc0016f66b7283ce0ce2c4055aa55
CPython 3.14.5
\end{lstlisting}
The root \path{lean-toolchain} and \path{lake-manifest.json} fix the Lean
and library revisions; \path{scripts/check-local-proof.py} selects the
Python version and records the checker identities used for local replay.

One \path{Challenge.lean}/\path{Solution.lean} pair represents both numbered
theorems through three declarations selected by \path{comparator.json}.
Within this pair, the challenge imports only Mathlib and contains explicit
mathematical definitions and the three statement holes.  Its definitions mirror
\path{DeepLean/Attention/Definitions.lean}, without importing project proof
modules.  The solution imports the completed library and proves declarations
with the names and types specified in the challenge.  Separate exports allow
Comparator to compare their statements without importing the challenge holes
into the solution.

For Theorem~\ref{thm:negative}, the declaration
\begin{lstlisting}
DeepLean.Attention.normalizer_counterexample
\end{lstlisting}
produces a scalar normalizer, proves its positivity,
normalization at zero, injectivity, Borel measurability, and discontinuity at
zero, establishes the
nonempty Euclidean-open dense locus $m_1m_2l\ne0$, and supplies at every point
the explicit distinct partner $(m_1,4m_2,l)$ with identical complete
realization. Its proof packages the explicit normalizer using the
admissibility, discontinuity, locus, and partner lemmas. The supporting endpoint
\begin{lstlisting}
Neuroalgebraic.Submission.Theorem1.theorem1_counterexample
\end{lstlisting}
states the existential counterexample. Discontinuity at zero for the explicit witness is also included in the
selected declaration.

For the pointwise clause of Theorem~\ref{thm:positive}, the declaration
\begin{lstlisting}
DeepLean.Attention.softmax_identifiability
\end{lstlisting}
requires a nonzero source final map, a non-skew proper source
score prefix, equal positive depths, and equality of the complete realization
functions at positive temperature with $t\ge2$.  No nondegeneracy condition
is imposed on the terminal source score or the competing parameter, and the
conclusion identifies the complete virtual parameter points.  Natural-number
dimensions are allowed formally; the
nonzero final map excludes zero-dimensional cases.  Its supporting endpoint
\begin{lstlisting}
Neuroalgebraic.Submission.Theorem2.theorem2_softmax_identifiability
\end{lstlisting}
has the identical pointwise statement.

For the explicit generic clause, the declaration
\begin{lstlisting}
DeepLean.Attention.softmax_generic_identifiability
\end{lstlisting}
states nonemptiness, Euclidean openness, and density of the coordinate locus
in Equation~\eqref{eq:coordinate-locus},
its containment in the source set in Equation~\eqref{eq:source-regular},
and singleton complete fibers throughout it.  A fixed-depth parameter with
tail length $k$ has depth $k+1$; at depth one the proper-prefix condition is
empty.  Existence of the coordinate indices witnesses positive input and output
dimensions, while the fiber ranges over all parameters of the fixed depth.
The generic wrapper combines the genericity and fiber proofs with
\path{NormalizedDeepAttention.properCoordinateRegular_properDeep}, making
the manuscript's containment clause explicit. Its supporting endpoint
\begin{lstlisting}
Neuroalgebraic.Submission.Theorem2.theorem2_openDense_coordinate_fibers
\end{lstlisting}
provides the openness, density, and singleton-fiber assertions.

The architecture reduction is supplied by
\begin{lstlisting}
NormalizedDeepAttention.deepNormalizedAttention_eq_virtual
\end{lstlisting}
which identifies the primitive and virtual realizations.

\noindent\begin{minipage}{\linewidth}
Log-affine reparametrization and boundary analysis are supplied by
\begin{lstlisting}
Neuroalgebraic.Submission.DeepAttentionVirtual.
  logAffine_properPrefix_pointwise_identifiable
  logAffine_properCoordinate_fullFiber
  properCoordinateLocus_nonempty_open_dense
  skewBoundary_nontrivialFiber
\end{lstlisting}
which remain supporting library results; Palomar comparison selects only
the three \path{DeepLean.Attention} declarations corresponding to the
numbered theorems.
\end{minipage}
\par

\subsection{Manuscript-facing Lean declarations}

Each theorem listing is imported directly from \path{Solution.lean}, so the
manuscript reproduces the source declaration without maintaining a separate copy.

\noindent\begin{minipage}{\linewidth}
\textbf{Theorem~\ref{thm:negative}.}\par
\smallskip
\lstinputlisting[
  xleftmargin=0.5em,
  xrightmargin=0.5em,
  linerange={APPENDIXTHEOREMONEBEGIN-APPENDIXTHEOREMONEEND},
  rangeprefix={--\ },
  includerangemarker=false
]{../../Solution.lean}
\end{minipage}
\par\medskip

\noindent\begin{minipage}{\linewidth}
\textbf{Theorem~\ref{thm:positive}: pointwise identifiability.}\par
\smallskip
\lstinputlisting[
  xleftmargin=0.5em,
  xrightmargin=0.5em,
  linerange={APPENDIXTHEOREMTWOPOINTWISEBEGIN-APPENDIXTHEOREMTWOPOINTWISEEND},
  rangeprefix={--\ },
  includerangemarker=false
]{../../Solution.lean}
\end{minipage}
\par\medskip

\newpage
\noindent\begin{minipage}{\linewidth}
\textbf{Theorem~\ref{thm:positive}: open-dense singleton fibers.}\par
\smallskip
\lstinputlisting[
  xleftmargin=0.5em,
  xrightmargin=0.5em,
  linerange={APPENDIXTHEOREMTWOGENERICBEGIN-APPENDIXTHEOREMTWOGENERICEND},
  rangeprefix={--\ },
  includerangemarker=false
]{../../Solution.lean}
\end{minipage}
\par\medskip

\subsection{Reproduction and evidence boundary}

With the pinned Lean toolchain installed through \path{elan} and the root
manifest's dependencies prepared, the complete local check is invoked from
the repository root by
\begin{lstlisting}
python3 scripts/check-local-proof.py
\end{lstlisting}
which creates a fresh evidence directory under \path{.cache/local-proof/}
and records source hashes, tool identities, commands, process exits, and a
final \path{report.json}.  The check validates the mirrored definitions and
declaration signatures, builds the project, elaborates Challenge and Solution
separately, scans completed sources for placeholders, and audits the transitive
axiom dependencies of all three selected declarations.  \path{Challenge.lean}
specifies the theorem statements with proof placeholders, while
\path{Solution.lean} supplies proofs containing no placeholders, directly or
through their dependencies.  The permitted axioms are exactly
\path{propext}, \path{Quot.sound}, and \path{Classical.choice}.

For proof replay, the script exports the selected declarations and required
primitive targets with \path{leanexport}.  Comparator compares the Challenge
and Solution statements and checks the axiom policy under
\path{comparator.json}, invoking NanoDa and Lean's checker; \path{con-ron},
an external Lean proof checker implemented in Rust, separately checks the
Solution export in verified mode.  Acceptance requires
results identifying every selected declaration and every required checker,
in addition to successful process exits.  

Each run records source checksums, tool identities, commands, and results in
its evidence directory under \path{.cache/local-proof/}.

\end{document}
